
\documentclass{book}

\usepackage{
  draftwatermark, % diagonal DRAFT watermark
  enumerate,      % enumerate with letter labels
  fancyhdr,       % custom chapter headings
  fullpage,       % reasonable margins
  hyperref,       % hotlinks
  makeidx,        % index
  tcolorbox,      % background colors in sidebars
  times,          % font choice
}

% Don't say "Chapter 3", just "3"
\renewcommand{\chaptername}{}

% do not indent paragraphs horizontally
\setlength{\parindent}{0pt}

% blank line between paragraphs
\setlength{\parskip}{12pt}

% draft watermark
\SetWatermarkText{DRAFT}
\SetWatermarkColor[gray]{0.9}

% tip sidebar
\newenvironment{tip}
    {\tcolorbox \begin{tabular}{p{2cm} l} \Large{TIP} & }
    {\end{tabular} \endtcolorbox}

\makeindex

\begin{document}

\title{Undergraduate Handbook \\ 2015-2016 Edition}
\author{Department of Computer Science \\ California State University, Fullerton}
\date{Draft: \today}
\maketitle

\newpage
\tableofcontents

\chapter{Introduction}

\section{The field of computer science}
Computer Science is the systematic study of computing systems and computation. The body of knowledge contains the theoretical foundation for understanding computing systems and methods, design methodology, algorithms, and software and hardware tools.

These programs cover a wide range of areas, including:
\begin{itemize}
\item multimedia and digital game technologies,
\item Internet and enterprise computing,
\item wireless and mobile computing,
\item databases and data mining,
\item computer security,
\item software engineering, and
\item computational bioinformatics.
\end{itemize}

Computer Science prepares graduates for rewarding careers in all areas of business, government, education and industry. These organizations, large and small, need computer professionals to address their needs with specific programs and systems. Computer science professionals tackle complicated problems and create computer solutions to solve them, devising new ways to use computers.

\section{The department}
The faculty and staff of the Computer Science Department welcome you into our program and sincerely wish you good luck on your journey into higher education, and continued success.

Whenever you have a question about the Department---its policies, its curriculum, its services, your progress, or anything else---feel free to contact us.

\index{contact information}
\begin{tabular}{rp{6in}} % right column in p mode so postal address can have newlines
  Web: & \href{http://fullerton.edu/ecs/cs/}{http://fullerton.edu/ecs/cs/} \\ \index{website}
  E-mail: & \href{mailto:csoffice@ecs.fullerton.edu}{csoffice@ecs.fullerton.edu} \\ \index{e-mail}
  In person: & Room CS-522 \\ \index{department office}
  Telephone: & (657) 278-3700 \\ \index{phone number} \index{telephone number}
  Fax: & (657) 278-7168 \\ \index{fax number}
  Postal mail: & California State University, Fullerton \newline \index{postal address} \index{address}
Department of Computer Science \newline
P.O. Box 6870 \newline
Fullerton, CA 92834-6870
\end{tabular}

\section{Accreditation} \index{accreditation}\index{ABET}

The Bachelor of Science in Computer Science degree at Cal State Fullerton is accredited by the Computer Accreditation Commission of ABET (\href{http://www.abet.org}{abet.org}).

\begin{center}
  \includegraphics{CAC-CMYK-W.pdf}
\end{center}

\section{The programs}

The Department offers the following Undergraduate programs, which are documented in this Handbook:
\begin{enumerate}
\item Bachelor of Science in Computer Science (BS CS), and
\item Minor in Computer Science.
\end{enumerate}

The Department also offers Graduate programs, which are documented elsewhere:
\begin{enumerate}
\item Master of Science in Computer Science (MS CS),
\item Master of Science in Software Engineering (MSE), and
\item Accelerated Master of Science in Software Engineering (AMSE).
\end{enumerate}

Some CS courses count toward the Computer Engineering, Electrical Engineering, and Mathematics majors.

\section{Objectives and Outcomes}

The following Program Educational Objectives and Program Outcomes have been established for the Bachelor of Science in Computer Science:

\textbf{Program Educational Objectives} \index{program educational objectives}
\begin{enumerate}
\item \emph{Technical Growth:}  Graduates will be successful in modern computing practices, integrate into the local and global workforce, and contribute to the economy of California and the nation.
\item \emph{Professional Skills:} Graduates will continue to demonstrate the professional skills necessary to be competent employees, assume leadership roles, and have career success and satisfaction.
\item \emph{Professional Attitude and Citizenship:} Graduates will become productive citizens with high ethical and professional standards, who make sound technical or managerial decisions, and have enthusiasm for the profession and professional growth.
\end{enumerate}

\textbf{Program Outcomes} \index{program outcomes}
\begin{enumerate}[(a)]
\item An ability to apply knowledge of computing and mathematics appropriate to the discipline
\item An ability to analyze a problem, and identify and define the computing requirements appropriate to its solution
\item An ability to design, implement, and evaluate a computer-based system, process, component, or program to meet desired needs
\item An ability to function effectively on teams to accomplish a common goal
\item An understanding of professional, ethical, legal, security and social issues and responsibilities
\item An ability to communicate effectively with a range of audiences
\item An ability to analyze the local and global impact of computing on individuals, organizations, and society
\item Recognition of the need for and an ability to engage in continuing professional development
\item An ability to use current techniques, skills, and tools necessary for computing practice
\item An ability to apply mathematical foundations, algorithmic principles, and computer science theory in the modeling and design of computer-based systems in a way that demonstrate comprehension of the tradeoffs involved in design choices
\item An ability to apply design and development principles in the construction of software systems of varying complexity
\end{enumerate}

\section{Using this document}

This handbook covers information on how to earn a Bachelor of Science or a Minor in Computer Science, and contains information relevant to students pursuing them. If you are pursuing a Masters degree, please trade this handbook for a copy of the Graduate Handbook.

Some aspects of our programs are complex, and you may find it difficult to choose among alternatives. In those cases, we present our suggested default choice as a tip, as shown below. You are not required to obey these tips, but obeying them is often a wise choice.
\begin{tip} \index{tips}
When in doubt, heed tips such as this one.
\end{tip}

\chapter{Other sources of information}

\chapter{The CS Major}

\section{Major requirements at a glance}

\section{Major prerequisite tree}

\section{Lower Division Core}
\section{Mathematics Core}
\section{Science and Mathematics Electives}
\section{Examination in Programming Proficiency (EPP)}
\section{Upper Division Core}
\section{Elective Tracks}
\section{Upper division writing requirement}
\section{General Education (GE)}
\section{Academic Requirements}
\section{Double-majoring in a related field}

\chapter{The CS Minor}

\section{Minor requirements at a glance}
\section{Minor prerequisite tree}
\section{Suggested minor electives}
\section{For majors in related fields}

\chapter{Major tracks}

\section{Multimedia and Digital Games (MG)}
\section{Internet and Enterprise Computing (IE)}
\section{Software Engineering (SE)}
\section{Scientific Computing (SC)}
\section{Customized (CT)}

\chapter{Alternative pathways}

\section{Transfer}

\section{Computer Science Placement Examination}

\section{Internships}

\section{Independent Study}

\chapter{Advisement}

\chapter{Resources and Activities}

\section{Open Labs}

\section{Tutoring Center}

\section{Supplemental Instruction}

\section{Clubs}
\subsection{Association of Computing Machinery (ACM)}
\index{ACM}

\subsection{ACM-W}
\index{ACM}\index{ACM-W}

\subsection{Upsilon Pi Epsilon (UPE)}
\subsection{Video Game Design Club (VGDC)}
\subsection{Security Club (?)}

\chapter{Progress Flowcharts}

\chapter{Credits and Revision History}

Copyright 2015, Department of Computer Science, California State University, Fullerton.

% magic to add index to table of contents without being a numbered chapter, from
% https://en.wikibooks.org/wiki/LaTeX/Indexing#Adding_index_to_table_of_contents
\cleardoublepage \addcontentsline{toc}{chapter}{Index} \printindex

\end{document}
